\chapter{Aspectos éticos y protección de datos}
\label{chap:gestion_etica}

\section{Responsabilidad clínica}

La posible aplicación como apoyo al triaje dejaría la revisión de las alertas y las decisiones diagnósticas y terapéuticas bajo responsabilidad del personal sanitario. Una predicción negativa no debería impedir una evaluación más exhaustiva si existe sospecha clínica, según la sección~\ref{sec:evaluacion_clinica}. Antes de incorporarla a la asistencia habría que evaluar los riesgos de omitir casos o generar falsas alertas y proteger también los datos que se recogieran durante ese proceso, con las limitaciones descritas en la sección~\ref{sec:limitaciones}.

La salida del modelo debe interpretarse junto con la historia clínica, la exploración y el resto de pruebas disponibles. En este TFG, la imagen del ECG y la etiqueta utilizada en la evaluación representan solo una parte del contexto necesario para valorar un caso. Por ello, una alerta puede servir para ordenar revisiones y facilitar la identificación de casos que merezcan atención, pero requiere una comprobación profesional antes de orientar cualquier decisión.

Si el sistema llegara a incorporarse a un circuito asistencial, habría que definir quién revisa cada alerta, qué información acompaña a la predicción y cómo se documenta la decisión final. También sería necesario vigilar cambios en la distribución de los registros y en la frecuencia de errores, porque el comportamiento observado en este conjunto puede variar al trabajar con otros centros, equipos o prácticas de adquisición. Estas condiciones deberían fijarse antes de su uso y revisarse durante el seguimiento.

\section{Protección de datos}

El artículo 9 del Reglamento General de Protección de Datos considera los datos de salud una categoría especial \cite{gdpr}. Su tratamiento exige una base jurídica adecuada, minimización, limitación de finalidad, confidencialidad y seguridad proporcional al riesgo. El artículo 89 añade salvaguardias para investigación científica.

En España, la Ley Orgánica 3/2018 incorpora garantías para investigación, en particular con datos seudonimizados \cite{lopdgdd}. La Ley 14/2007 añade principios de confidencialidad, supervisión ética y protección de participantes \cite{ley142007}.

En este trabajo, el conjunto cedido por el HUCA declara la anonimización mediante Safe Harbor \cite{costacortez2026brugada}. El análisis utiliza únicamente los campos necesarios para construir las imágenes, formar los pliegues y evaluar las predicciones. Los identificadores directos quedan fuera de los materiales de difusión y los resultados se presentan de forma agregada, de manera que las tablas y las figuras no permiten reconstruir la trayectoria de un paciente concreto.

La protección debe mantenerse durante todo el ciclo del estudio. Esto incluye la preparación de las imágenes, la generación de los mensajes para Gemma~4, el almacenamiento de las predicciones y la elaboración de los resultados. La reutilización de una imagen junto con información clínica adicional podría aumentar el riesgo de identificación por combinación de variables, incluso cuando cada campo aislado pareciera inocuo. Por ese motivo, cualquier unión con otras fuentes debe quedar fuera del flujo ordinario del experimento hasta revisar sus garantías y su finalidad.

La elección del entorno de ejecución también forma parte de esta protección. Las imágenes permanecen en infraestructura universitaria autorizada y la inferencia de Gemma~4 se realiza mediante un servicio controlado en ese entorno. Esta configuración reduce la exposición a servicios externos, aunque no elimina las obligaciones asociadas al tratamiento de datos de salud. El acceso debe limitarse a las personas que participan en el proyecto, conservarse durante el tiempo imprescindible y quedar sujeto a medidas de autenticación, autorización y registro. En el caso de los modelos de visión y lenguaje, estas precauciones deben aplicarse tanto a las imágenes incluidas en los mensajes como a las respuestas y a los registros técnicos que pueda generar el servicio.

El carácter anonimizado del conjunto facilita la investigación, pero no convierte cualquier uso posterior en automáticamente admisible. Una ampliación de las variables, una vinculación con otras fuentes o una reutilización con finalidad asistencial exigiría comprobar de nuevo la base jurídica, la finalidad, los plazos de conservación y las garantías aplicables. También sería necesario valorar una evaluación de impacto cuando el tratamiento o el riesgo cambiasen de forma relevante. Antes de trasladar el sistema al triaje, estas decisiones deberían acompañarse de supervisión ética, control de acceso y validación clínica, manteniendo la revisión profesional de cada caso.
